\usepackage[a4paper, margin=1in]{geometry}
\usepackage{minted}
\setminted{breaklines, % Break lines automatically
    breakanywhere, % Allow breaks anywhere
    fontsize=\small, % Font size for the code
    frame=single, % Border style for the code block
    bgcolor=darkbackground, % Background color for the code block
    linenos, % Display line numbers
    numbersep=5pt, % Space between line numbers and code
    label=code-\thecodeblock}
\usepackage{caption}
\usepackage{tcolorbox}
\tcbuselibrary{breakable, skins, theorems, listings}
\usepackage{titlesec}
\usepackage{fancyhdr}
\usepackage{pgf-pie}
\usepackage{filecontents}
\usepackage{hyperref}
\usepackage{lipsum}
\usepackage{amsmath}
\usepackage{amsthm}
\usepackage{amssymb}
\usepackage{enumitem}
\usepackage{tikz}
\usepackage{tocloft}
\usepackage{xparse}
\usepackage{chngcntr}
\usepackage{mathtools}
\usepackage{thmtools} 
\usepackage{pgfplots}
\usepackage{siunitx}
\usepackage{chemfig}
% For flexible counter manipulation

% Redefine equation numbering to chapter.section.subsection.number
\counterwithin{equation}{subsection}

\newenvironment{introduction}{
    \chapter*{Introduction}
    \addcontentsline{toc}{chapter}{Introduction}
}{
    % You can put any ending code here if necessary
}
\newenvironment{acknowledgements}{
    \chapter*{Acknowledgements}
    \addcontentsline{toc}{chapter}{Acknowledgements}
}{
    % You can put any ending code here if necessary
}
\newenvironment{references}{
    \chapter*{References}
    \addcontentsline{toc}{chapter}{References}
    \begin{thebibliography}{99} % Start bibliography
}{
    \end{thebibliography} % End bibliography
}

% Ensure equation counter resets with each new subsection
\renewcommand{\theequation}{\thesubsection.\arabic{equation}}

% Colors
\usepackage{xcolor}
\definecolor{whitetext}{HTML}{121212}
\definecolor{darkbackground}{HTML}{FFFFFF}
\definecolor{accent3}{HTML}{FF5C00}
\definecolor{lightcyan}{HTML}{E0FFFF}
\definecolor{lightgray}{HTML}{D3D3D3}
\definecolor{accent1}{HTML}{0047AB}
\definecolor{accent3}{HTML}{6495ED}
\definecolor{accent3}{HTML}{FF4D00}
\definecolor{myorange}{HTML}{FFA500}
\definecolor{mylightblue}{HTML}{ADD8E6}
\definecolor{mygreen}{HTML}{90EE90}
\definecolor{definition}{HTML}{924492}
\definecolor{axiom}{HTML}{E23458}

\pagecolor{darkbackground}
\color{whitetext}

\newenvironment{longlisting}{\captionsetup{type=listing}}{}

% Custom Counters
\newcounter{mydefinition}[subsection]
\renewcommand{\themydefinition}{\thesection.\thesubsection.\thesubsubsection.\arabic{mydefinition}}
\newcounter{myidentity}[subsection]
\renewcommand{\themyidentity}{\thesubsection.\arabic{myidentity}}
\newcounter{myproof}[subsection]
\renewcommand{\themyproof}{\thesubsection.\arabic{myproof}}
\newcounter{mynote}[subsection]
\renewcommand{\themynote}{\thesubsection.\arabic{mynote}}
\newcounter{mycodeblock}[subsection]
\renewcommand{\themycodeblock}{\thesubsection.\arabic{mycodeblock}}
\newcounter{myexercise}[subsection]
\renewcommand{\themyexercise}{\thesubsection.\arabic{myexercise}}
\newcounter{codeblock}[section]
\renewcommand{\thecodeblock}{\thesubsection.\arabic{codeblock}}
\newcounter{myaxiom}[subsection]
\renewcommand{\themyaxiom}{\thesubsection.\arabic{myaxiom}}


\setminted{breaklines, % Break lines automatically
    breakanywhere, % Allow breaks anywhere
    fontsize=\small, % Font size for the code
    frame=single, % Border style for the code block
    bgcolor=darkbackground, % Background color for the code block
    linenos, % Display line numbers
    numbersep=5pt, % Space between line numbers and code
    label=code-\thecodeblock}

% Custom Environments
\tcbuselibrary{listingsutf8}
\newtcolorbox[use counter=mydefinition, number within=subsection]{mydefinition}[2][]{colback=darkbackground, colframe=definition, fonttitle=\bfseries, breakable=true, enhanced jigsaw, sharp corners, parbox=false, title=Definition~\themydefinition, before upper={\textbf{#2}\\}, #1}
\newtcolorbox[use counter=myidentity, number within=subsection]{myidentity}[1][]{colback=darkbackground, colframe=myorange, fonttitle=\bfseries, breakable=true, enhanced jigsaw, sharp corners, before skip=10pt, after skip=10pt, parbox=false, title=Identity~\themyidentity, #1}
\newtcolorbox[use counter=myproof, number within=subsection]{myproof}[1][]{colback=darkbackground, colframe=mygreen, fonttitle=\bfseries, breakable=true, enhanced jigsaw, sharp corners, before skip=10pt, after skip=10pt, parbox=false, title=Proof~\themyproof, #1}
\newtcolorbox[use counter=mynote, number within=subsection]{mynote}[1][]{colback=darkbackground, colframe=accent3, fonttitle=\bfseries, breakable=true, enhanced jigsaw, sharp corners, before skip=10pt, after skip=10pt, parbox=false, title=Note~\themynote, #1}
\newtcolorbox[use counter=mycodeblock, number within=subsection]{mycodeblock}[1][]{colback=darkbackground, colframe=accent3, fonttitle=\bfseries, breakable=true, enhanced jigsaw, sharp corners, before skip=10pt, after skip=10pt, parbox=false, title=Code~\themycodeblock, #1}
\newtcolorbox[use counter=myexercise, number within=subsection]{myexercise}[1][]{colback=darkbackground, colframe=accent1, fonttitle=\bfseries, breakable=true, enhanced jigsaw, sharp corners, before skip=10pt, after skip=10pt, parbox=false, title=Exercise~\themyexercise, #1}
\newtcolorbox[use counter=myaxiom, number within=subsection]{myaxiom}[1][]{colback=darkbackground, colframe=axiom, fonttitle=\bfseries, breakable=true, enhanced jigsaw, sharp corners, before skip=10pt, after skip=10pt, parbox=false, title=Axioms ~\themyaxiom, #1}
% Custom Theorem-like Environments
\newtheorem{mytheorem}{Theorem}[chapter]
\newtheorem{mydefinitiontheorem}{Definition}[subsection]
\newtheorem{myexercisetheorem}{Exercise}[subsection]

% Syntax Highlighting
\usemintedstyle{monokai}
\setminted{
    bgcolor=darkbackground,
    style=monokai,
    linenos,
    breaklines,
    fontsize=\small
}

% Custom Title Page with TikZ Decorations and Logo
\newcommand{\customtitle}[6]{
    \begin{titlepage}
        \noindent
        \begin{tikzpicture}[remember picture,overlay]
            \begin{scope}[blend mode=soft light]
                \fill[accent1] (current page.south west) rectangle (current page.north east);
                \fill[accent3] (0.2*\paperwidth,0.2*\paperheight) -- (current page.north east) -- (current page.north west) -- cycle;
                \fill[accent3] (current page.south west) -- (current page.north west) -- (0.8*\paperwidth,0.8*\paperheight) -- cycle;
                \draw[line width=1mm, accent3, opacity=0.5] (current page.south west) -- (current page.north east);
                \draw[line width=1mm, accent3, opacity=0.5] (current page.south east) -- (current page.north west);
            \end{scope}
            % Logo
            \node [anchor=north, yshift=-8.2cm] at (current page.north) {
                \includegraphics[width=0.5\textwidth]{#6}
            };
            % Title and Subtitle in cyan box
            \node [fill=accent3, fill opacity=0.8, text opacity=1, inner sep=1cm, anchor=north, text width=\paperwidth, minimum height=0.15\paperheight, align=center] at ([yshift=-2cm]current page.center) {
                \Huge\textbf{\textit{#1}}\\[0.5cm]
                \Large\textbf{\textit{#5}}
            };
            % Author, Subject, and Date in royal blue box
            \node [fill=accent3, fill opacity=0.8, text opacity=1, inner sep=0.5cm, anchor=north, text width=\textwidth-8cm, align=center] at ([yshift=-6.45cm]current page.center) {
                \Large\textbf{\textit{#2}}\\[0.5cm]
                \normalsize\textbf{#3}\\[0.5cm]
                \normalsize\textbf{#4}
            };
        \end{tikzpicture}
    \end{titlepage}
}

% Custom Chapter Page with TikZ Decorations
\newcommand{\chapterpage}{
    \begin{tikzpicture}[remember picture,overlay]
        \begin{scope}[blend mode=soft light]
            \fill[accent1] (0,0) rectangle (\paperwidth,\paperheight);
            \fill[accent3] (0.2*\paperwidth,0.2*\paperheight) -- (\paperwidth,0.8*\paperheight) -- (\paperwidth,0) -- cycle;
            \fill[accent3] (0,0) -- (\paperwidth,0) -- (0.8*\paperwidth,0.8*\paperheight) -- cycle;
        \end{scope}
    \end{tikzpicture}
}

%Custom math commands
\newcommand{\intg}[3]{\int_{#1}^{#2} #3 \, dx}
\newcommand{\sumg}[3]{\sum_{#1}^{#2} #3}
\newcommand{\diff}[2]{\frac{d #1}{d #2}}
\newcommand{\pdiff}[2]{\frac{\partial #1}{\partial #2}} % Partial differential
\newcommand{\limit}[3]{\lim_{#1 \to #2} #3} % Limit
\newcommand{\infseries}[2]{\sum_{#1}^{\infty} #2} % Infinite series
\newcommand{\mat}[1]{\begin{matrix} #1 \end{matrix}} % Matrix
\newcommand{\pmat}[1]{\begin{pmatrix} #1 \end{pmatrix}} % Parentheses matrix
\newcommand{\bmat}[1]{\begin{bmatrix} #1 \end{bmatrix}} % Bracket matrix
\newcommand{\mycases}[1]{\left\{\begin{array}{lr}
        x(n), & \text{for } 0\leq n\leq 1\\
        x(n-1), & \text{for } 0\leq n\leq 1\\
        x(n-1), & \text{for } 0\leq n\leq 1
        \end{array}\right\}}
\newcommand{\floor}[1]{\left\lfloor #1 \right\rfloor} % Floor function
\newcommand{\ceil}[1]{\left\lceil #1 \right\rceil} % Ceiling function
\newcommand{\abs}[1]{\left| #1 \right|} % Absolute value
\newcommand{\norm}[1]{\left\| #1 \right\|} % Norm
\newcommand{\set}[1]{\left\{ #1 \right\}} % Set
\newcommand{\angg}[1]{\left\langle #1 \right\rangle} % Angle
\newcommand{\eqn}[1]{\begin{equation} #1 \end{equation}} % Equation
\newcommand{\eqnarrayx}[1]{\begin{eqnarray} #1 \end{eqnarray}} % Equation array
\newcommand{\R}{\mathbb{R}}
\newcommand{\Z}{\mathbb{Z}}
\newcommand{\N}{\mathbb{N}}
\newcommand{\Q}{\mathbb{Q}}
\newcommand{\forallx}{\forall x \in \mathbb{R}}
\newcommand{\existsx}{\exists x \in \mathbb{R}}
\newcommand{\iintd}[1]{\iint\!#1\,dx\,dy}
\newcommand{\iiintd}[1]{\iiint\!#1\,dx\,dy\,dz}
\newcommand{\nint}[3]{\int_{#1}^{#2}\!#3\,dx}
\newcommand{\C}{\mathcal{C}}
\newcommand{\nd}{\wedge}
\newcommand{\orr}{\vee}
\newcommand{\union}{\cup}
\newcommand{\intersect}{\cap}
\newcommand{\megaunion}[2]{\bigcup\limits_{#1}^{#2}}
\newcommand{\megaintersect}[2]{\bigcap\limits_{#1}^{#2}}
\newcommand{\emphasise}[1]{\textbf{\textit{#1}}} % Emphasize text
\newcommand{\underemphasise}[1]{\underline{\textbf{\textit{#1}}}} % Custom title formatting

% Custom Headers and Footers
\pagestyle{fancy}
\fancyhf{}
\fancyhead[LE,RO]{\thepage}
\fancyhead[RE]{\leftmark}
\fancyhead[LO]{\rightmark}

% Section Formatting
\titleformat{\section}
  {\normalfont\Large\bfseries\color{accent3}}
  {\thesection}{1em}{}
\titleformat{\subsection}
  {\normalfont\large\bfseries\color{accent3}}
  {\thesubsection}{1em}{}
\titleformat{\subsubsection}
  {\normalfont\normalsize\bfseries\color{accent3}}
  {\thesubsubsection}{1em}{}

\newcommand{\settitle}[1]{
\hypersetup{
    colorlinks=true,
    linkcolor=accent3,
    filecolor=accent3,
    urlcolor=accent3,
    citecolor=accent3,
    pdftitle={#1},
    pdfpagemode=FullScreen
}}
\newcounter{example}[section] % Reset example counter with each section
\renewcommand{\theexample}{\thesubsection.\arabic{example}} % Customize numbering

\newenvironment{example}[1][]{%
    \refstepcounter{example}% Increment example counter
    \noindent\textit{\textbf{Example~\theexample~ - #1}}
    \par\nobreak % Print the title and ensure no page break% Add vertical space
    \noindent % No indentation for the first paragraph
}{% % Ensure a new paragraph ends the environment
    \par 
    \noindent
}

\newcounter{solution}[chapter] % Reset example counter with each section
\renewcommand{\thesolution}{\thechapter.\arabic{solution}} % Customize numbering

\newenvironment{solution}[1][]{%
\refstepcounter{solution}
\textit{\textbf{Solutions~\thesolution:}} \textit{#1}\par\nobreak
\vspace{0.5\baselineskip}
\noindent}{%
\par \vspace{0.9\baselineskip} \noindent}
\newcounter{exercise}[chapter] % Reset exercise counter with each section
\renewcommand{\theexercise}{\thechapter.\arabic{exercise}} % Customize numbering
\newenvironment{exercise}[1][]{%
    \refstepcounter{exercise}% Increment exercise counter
    \noindent\underline{\textbf{Exercise~\theexercise:}}
}{% % Ensure a new paragraph ends the environment
    }
\newcounter{corollary}[mytheorem] % Reset corollary counter with each theorem
\newenvironment{corollary}[1][]{%
    \refstepcounter{corollary}
    \noindent\underline{\textbf{Corollary~\themytheorem.\thecorollary:}} \textit{#1}\par\nobreak
    \vspace{0.5\baselineskip}
    \noindent
}{
}